\chapter{测试矩阵与 CI}
\label{ch:test-ci}

\section{测试规模概览（\versiontag{0.9.4}）}

\begin{itemize}[nosep]
  \item \textbf{232} GoogleTest \texttt{TEST()} 宏（\texttt{ebbackup\_tests} target）；
  \item \textbf{76} 测试源文件（\filepath{*\_test.cc}）；
  \item C API：\texttt{eb\_run\_tests\_capi}（8 tests）；
  \item Bench 门禁：\texttt{ebbackup\_bench\_check}（L1--L7）；
  \item CI：\filepath{.github/workflows/ebbackup.yml} — Windows + Linux Release + Linux ASan。
\end{itemize}

\section{232 测试按目录分解}

\begin{longtable}{@{}rlrp{6.5cm}@{}}
\toprule
目录 & Tests & 文件数 & 覆盖重点 \\
\midrule
\endhead
\filepath{tests/engine/} & 79 & 24 & E2E、restore、snapshot、manifest、C API、fixture \\
\filepath{tests/store/} & 23 & 11 & ChunkStore、EbPack、index、compact、GC、repo-stats \\
\filepath{tests/failure/} & 20 & 4 & powerfail、chaos、balanced durability \\
\filepath{tests/chunk/} & 20 & 7 & FastCDC、streaming、digest、HCRBO bounds \\
\filepath{tests/pipeline/} & 20 & 7 & v3/v4、streaming、parity、encrypt \\
\filepath{tests/crypto/} & 21 & 6 & AES-GCM NIST、PBKDF2、digest 双栈 \\
\filepath{tests/scan/} & 16 & 3 & backup filter（12）、special files \\
\filepath{tests/audit/} & 8 & 3 & RAR、Merkle、CARL \\
\filepath{tests/daemon/} & 9 & 4 & schedule JSON、watch smoke \\
\filepath{tests/codec/} & 8 & 3 & LZ4、zstd、ContentClass \\
\filepath{tests/state/} & 5 & 3 & superblock、FSM、sync executor \\
\filepath{tests/cli/} & 3 & 1 & restore CLI \\
\filepath{tests/archive/} & 1 & 1 & eb\_bundle \\
\filepath{tests/bench/} & 1 & 1 & reuse gate \\
\midrule
\textbf{合计} & \textbf{232} & \textbf{76} & \\
\bottomrule
\end{longtable}

\section{Powerfail：EBTEST\_ABORT\_AFTER}

\texttt{EBTEST\_ABORT\_AFTER=<int>} 设为 \texttt{BackupPhase} 枚举整数值，在对应 phase \texttt{MaybeTestAbortAfter} 返回 \texttt{Internal("test abort after phase")}，模拟 crash。

\begin{table}[htbp]
\centering
\caption{BackupPhase 注入值对照}
\begin{tabular}{@{}ll@{}}
\toprule
Phase & 枚举值（十进制） \\
\midrule
\texttt{kIdle} & 0 \\
\texttt{kScanning} & 1 \\
\texttt{kChunking} & 2 \\
\texttt{kStoring} & 3 \\
\texttt{kCommittingMeta} & 4 \\
\texttt{kAuditing} & 5 \\
\texttt{kComplete} & 6 \\
\texttt{kAborted} & 255 \\
\bottomrule
\end{tabular}
\end{table}

测试套件：
\begin{itemize}[nosep]
  \item \texttt{powerfail\_test.cc}（4）：双 slot corrupt、interrupt recovery、subprocess kill；
  \item \texttt{powerfail\_extended\_test.cc}（8）：pipeline strict/balanced 多 kill 点；
  \item 子进程模式：\texttt{RunBackupSubprocessAndKill} 对 CLI \texttt{eb backup} 发 signal。
\end{itemize}

验证断言：txn $\rightarrow$ \texttt{kAborted}；前一 snapshot \texttt{verify} OK；Recover 后可重新 backup。

\section{Chaos 测试（seed=42）}

\filepath{tests/failure/chaos\_test.cc}（4 tests）：

\begin{itemize}[nosep]
  \item \texttt{RandomPhaseThenRecover}：\texttt{std::mt19937 gen(42)}；8 trials（\texttt{kChaosTrials=8}）；随机 phase abort + Recover；
  \item \texttt{RandomChunkBitFlipThenVerify}：chunk 字节翻转 $\rightarrow$ verify corrupt；
  \item \texttt{chaos\_util}：封装 \texttt{ScopedAbortAfterPhase}、子进程 kill；
  \item v0.4.2 起固定 seed=42 保证 CI 确定性。
\end{itemize}

\section{Fixture 清单}

\begin{longtable}{@{}llp{7.5cm}@{}}
\toprule
Fixture & 路径/变量 & 内容 \\
\midrule
Synthetic & \texttt{MakeSyntheticData(size, seed)} & 确定性伪随机字节 \\
Real-world & \texttt{EBTEST\_FIXTURE\_DIR}、\texttt{real\_fixture\_test} & 5 层嵌套目录、Unicode 文件名、多类型 \\
Media & \texttt{tests/fixtures/media/}、\texttt{media\_fixture\_test} & JPEG/WebP/MP4/ZIP 等 $\sim$2MB checked-in \\
Temp root & \texttt{EBTEST\_TMPDIR} & 默认 \filepath{test\_output/} \\
Engine root & \texttt{EBTEST\_ENGINE\_ROOT} & 子进程 CLI 查找 \texttt{eb} 二进制 \\
CI 标记 & \texttt{EBTEST\_CI=1} & 缩减随机 trial（如 powerfail 16$\rightarrow$4） \\
\bottomrule
\end{longtable}

Init helpers：\texttt{InitLegacyRepo}、\texttt{InitV03Repo}、\texttt{InitV05Repo}/\texttt{InitDefaultRepo}（见 \cref{ch:architecture}）。

\section{CI Workflow 步骤}

\filepath{.github/workflows/ebbackup.yml}：

\begin{enumerate}
  \item \textbf{Trigger}：\filepath{engine\_cpp/**}、\filepath{CMakeLists.txt}、\filepath{gtest\_capi/**} push/PR；
  \item \textbf{test-windows}：\texttt{cmake -DEBBACKUP\_BUILD\_TESTS=ON} $\rightarrow$ Release build $\rightarrow$ \texttt{ctest -C Release}（含 bench gate）；
  \item \textbf{test-linux}：同上，并行编译 \texttt{-j} 等于 CPU 核数；
  \item \textbf{test-linux-asan}：\texttt{-fsanitize=address,undefined} 构建 $\rightarrow$ \texttt{ctest}（无 Release 配置）。
\end{enumerate}

所有 job 设置 \texttt{EBTEST\_CI=1}。

\section{ASan Job 要点}

\begin{itemize}[nosep]
  \item 检测 ChunkStore 并发 rehash、use-after-free（v0.9.0 修复的 pipeline store race）；
  \item \texttt{-fno-omit-frame-pointer -g} 保留栈回溯；
  \item 与 Release bench 互补：ASan 不重跑 L5 绝对吞吐门禁（太慢），但跑全量 232 gtest。
\end{itemize}

\section{本地等效命令}

\begin{lstlisting}[language=bash,caption={本地 CI 等效}]
cmake -S . -B build -DEBBACKUP_BUILD_TESTS=ON
cmake --build build --config Release
set EBTEST_CI=1
ctest --test-dir build -C Release --output-on-failure
\end{lstlisting}

\section{本章小结}

232 gtest + L1--L7 bench + 双平台 CI + ASan 构成 ebbackup 交付质量网。附录提供完整源码索引与环境变量表。
